\section{Final Specifications and Requirements}
The system was designed and evaluated under the following specifications.

\textbf{Hardware and compute.} All per-modality training and fusion training was executed on Google Colab Pro using an NVIDIA L4 GPU with 24 GB VRAM and 83 TFLOPS FP16 performance. Local development and analysis were performed on a Windows PC with Intel i7 CPU and 32 GB RAM, and in earlier phases on an Apple M1 MacBook Air. The absence of local GPU compute necessitated a hybrid workflow in which code is authored locally and executed remotely, with all data and checkpoints persisted to Google Drive at \texttt{MyDrive/Thesis\_EAV/}.

\textbf{Data specifications.} The EAV dataset provides pickle files of pre-extracted features for each of 42 subjects across three modalities. Each audio pickle contains Log-Mel Spectrogram segments of shape \texttt{(400, time, 128)} per subject, resampled to 16 kHz and segmented into five-second clips. Each vision pickle contains 25 video frames per clip, optionally face-cropped to $224\times224$ pixels. Each EEG pickle contains five-second signal windows at 100 Hz after downsampling from 500 Hz, band-pass filtered to 0.5--45 Hz, across 30 channels. The train/test split uses a stratified-by-class scheme (\texttt{h\_idx=56}) yielding 280 train and 120 test trials per subject per modality.

\textbf{Software and libraries.} The implementation uses PyTorch 2.10, the Hugging Face Transformers library 4.4x, torchaudio, torchvision, librosa, NumPy, Pandas, and scikit-learn. The EAV repository provides base data loaders and per-modality trainers, modified to fix the bugs documented in Chapter~\ref{chap:engineering_challenges}. The fusion module is implemented in a standalone \texttt{fusion/} directory at the project root, with no dependency on EAV repository internals.

\textbf{Training specifications.} Per-modality epoch budgets are: AST audio, 5 frozen + 5 fine-tuning epochs; ViT vision, 3 frozen + 2 fine-tuning epochs; EEGNet, 350 epochs from scratch. Within-subject fusion training uses 80 epochs, the AdamW optimiser, learning rate $10^{-3}$, cosine annealing and batch size 32, and the final epoch is evaluated with no epoch selection. Cross-subject fusion heads use 40 epochs at batch size 64 with the best epoch chosen on three inner-validation participants. The full 42-subject within-subject rollout required approximately 12 hours of wall-clock time on the Colab L4 GPU; the cross-subject evaluation retrains all three encoders for each of five folds (Chapter~\ref{chap:crosssubject}).

\textbf{Code and reproducibility.} All code is publicly archived in the GitHub repository \texttt{shafinmuffinn/thesis\_p2}. A single path-configuration module, \texttt{paths.py}, reads all filesystem locations from environment variables, enabling the same code to run unchanged in both local and Colab environments. Per-subject, per-modality logits are saved as compressed NumPy archives, enabling all downstream fusion and coherence analysis to be reproduced from cached results without re-running the expensive per-modality training.

\section{Societal and Clinical Impact}
The societal impact of this research runs through three channels.

\textbf{Emotion recognition that works for new users.} Any deployed emotion recognition system meets people it was not trained on. The subject-independent evaluation of Chapter~\ref{chap:crosssubject} measures how much accuracy is lost in that situation on EAV, and the calibration method of Chapter~\ref{chap:calibration} recovers most of it with a short unlabelled enrolment recording and no labels. For applications such as adaptive human-computer interfaces, this is the difference between a system that works on its developers and one that works on its users.

\textbf{More reliable research practice.} The audit of Chapter~\ref{chap:audit} shows how easily a plausible result on this benchmark can be produced by test-set model selection, and how easily a map of one classifier's errors can be read as a map of human emotion regulation. The tests it uses are cheap and reusable, and they apply to every published fusion result on EAV.

\textbf{Mental health monitoring, with caution.} Detecting divergence between expressed and felt emotion remains a clinically motivated goal, relevant to conditions such as alexithymia and depression in which self-report is unreliable. This thesis does not provide such a tool. Its contribution to that goal is negative but useful: it shows that classifier disagreement alone is not evidence of suppression, and specifies the permutation test and the independent internal-state measure that a future clinical study would need before any claim could be made.

\section{Ethical Considerations}
The EAV dataset is an ethically approved corpus collected by Lee, Shomanov, Kabidenova, Yazici, and colleagues at Nazarbayev University. Participants provided informed consent to recording and to the release of de-identified data for research purposes. The dataset is not openly distributed; access is granted on individual request to the dataset authors. For this thesis, access was granted by Adai Shomanov via a token-authorised Zenodo link, with explicit permission for the dataset's use in this research. This thesis uses the pre-processed EAV pickle files in exactly the conditions intended by the data providers; no additional data collection or human-participants research was conducted.

Two further considerations follow from the results. First, the features learned for emotion recognition identify the individual: a linear classifier recovers which of the 42 participants produced a trial with 97--100\% accuracy from audio or video features and 82\% from EEG features (Section~\ref{sec:cal_mechanism}). Stored emotion features are therefore personal data, and any deployment must handle them as such; the calibration method additionally requires an enrolment recording of each user, which needs informed consent and secure handling. Second, coherence or suppression analyses applied to clinical populations raise considerations of privacy, stigmatisation, and clinical deployment without appropriate validation---all the more so given that Chapter~\ref{chap:audit} finds the present suppression matrix to be indistinguishable from classifier error. These are explicitly out of scope for the present thesis, which is confined to analysis of laboratory data. Any future work involving clinical populations must obtain independent ethics approval and must not make clinical inferences based solely on the present thesis's findings, which have not been validated against clinical ground truth.
